\documentclass[11pt]{article}
\usepackage{mathblatt}
% gesetzt von werkzeuge/setzer.py aus dem Lernweg (L1-A · L1-B · L1-C · L1-D · L1-T) und den Bankzeilen; Muster T6B.
% Präambel des Setzers (werkzeuge/setzer.py), wortgleich aus dem Hand-Satz T6B (bau/proben/2026-10-09/kathete-berechnen), dort aus K4W.
\makeatletter
\setlength{\mbpfbe}{0mm}% keine Punkte (Bauregel 6.4)
% eingerückt (Bauregel 4.7, 6.6): \gEin Einzug, \gSchrift eine Stufe kleiner
\newlength{\gEin}\setlength{\gEin}{0pt}
\let\gSchrift\relax
\renewcommand{\pfkreuzzeile}[1]{\par\noindent\hangindent3.2em\hangafter1\hspace*{1.6em}$\square$\ #1\par}
\newcommand{\gkreuz}[1]{$\square$\ #1\hspace{2em}}
% Bauregel 6.7: links, was man ansieht, rechts, was man tut; Spaltengrenze überall an derselben Stelle
\newsavebox{\gbox}\newlength{\gLinks}
\newcommand{\gzwei}[2]{\par\smallskip\noindent\sbox{\gbox}{#1}%
  \setlength{\gLinks}{\dimexpr0.5\textwidth-\gEin-\mbpfnr\relax}%
  \ifdim\wd\gbox>\dimexpr\gLinks-4mm\relax
    \usebox{\gbox}\par\smallskip\noindent\begin{minipage}{\linewidth}#2\end{minipage}%
  \else
    \begin{minipage}[t]{\gLinks}\vspace{0pt}\usebox{\gbox}\end{minipage}%
    \begin{minipage}[t]{\dimexpr\linewidth-\gLinks\relax}\vspace{0pt}\raggedright #2\end{minipage}%
  \fi\par}
\RenewDocumentEnvironment{pfaufg}{m m m}{%
  \begin{lrbox}{\mbaufgabebox}\begin{minipage}[t]{\linewidth}%
  \gSchrift\noindent\hspace*{\gEin}\mbpfrandmarke{#1}%
  \begin{minipage}[t]{\mbpfnr}\strut\textbf{#2}\end{minipage}%
  \begin{minipage}[t]{\dimexpr\linewidth-\gEin-\mbpfnr\relax}\raggedright\strut\ignorespaces}%
 {\end{minipage}%
  \end{minipage}\end{lrbox}%
  \setlength{\mbaufgabehoehe}{\dimexpr\ht\mbaufgabebox+\dp\mbaufgabebox\relax}%
  \par\addvspace{7pt}\Needspace*{\mbaufgabehoehe}%
  \noindent\usebox{\mbaufgabebox}\par\addvspace{7pt}}
\makeatother
% Rechenraum als Karo (Bauregel 6.9): n Rechenschritte = n mal 10 mm
\newcommand{\karo}[1]{\par\smallskip\noindent\begin{tikzpicture}
  \clip (0,0) rectangle ({\linewidth-1mm},{#1*10mm});
  \draw[black!16,line width=0.3pt,step=5mm] (0,0) grid ({\linewidth},{#1*10mm});
  \end{tikzpicture}\par}
\newcommand{\kk}[1]{$\square$\,#1\hspace{0.9em}}
\newcommand{\linie}[1]{\,{\color{black!45}\rule[-1pt]{#1}{0.4pt}}\,}
% Zeile am Ende jedes Durchgangs: Originale klein grau, darüber; dann Vorgänger/Nachfolger (Bauregel 3.4, 6.5)
\newcommand{\blattende}[2]{\par\vfill\noindent\begin{minipage}{\linewidth}{\scriptsize\color{mbgrau}Originale: #1\par}\smallskip
{\footnotesize #2\par}\end{minipage}}
\newcommand{\pkt}{\textperiodcentered{}}

\newcommand{\kennung}{VUD}
\newcommand{\grau}[1]{{\color{mbgrau}#1}}

\begin{document}
\pfheftstil
\fancyfoot[C]{\parbox[t]{\textwidth}{\mbfhbox\par\vspace{1.5mm}{\footnotesize\color{mbgrau}{\scriptsize \kennung}\hfill\thepage}}}
\pfheftkopf{Potenzen}{Klasse 9}
% L1-A: Potenz als Malkette
\begin{pfaufg}{}{1.}{}
\grau{a)\ Schreibe $6^3$ als Malkette und rechne aus. Schreibe $5 \cdot 5 \cdot 5 \cdot 5$ als Potenz und rechne aus.\par\smallskip Unten steht die Basis $6$, oben der Exponent $3$. Der Exponent sagt: drei Faktoren $6$. \quad $6^3 = 6 \cdot 6 \cdot 6$\par Schritt für Schritt malnehmen: \quad $6 \cdot 6 = 36$, \quad $36 \cdot 6 = 216$. \ Also $6^3 = 216$.\par Rückwärts: In $5 \cdot 5 \cdot 5 \cdot 5$ steht viermal der Faktor $5$. Basis $5$, Exponent $4$: \quad $5^4$\par Ausrechnen: \quad $5 \cdot 5 = 25$, \ $25 \cdot 5 = 125$, \ $125 \cdot 5 = 625$. \ Also $5^4 = 625$.}
\par\medskip
b)\ Schreibe als Malkette und rechne aus: \ a)~$2^5$ \qquad b)~$10^4$ \\ Schreibe als Potenz und rechne aus: \ c)~$7 \cdot 7 \cdot 7$ \qquad d)~$3 \cdot 3 \cdot 3 \cdot 3 \cdot 3$\karo{3}
\fusshilfe{\mbox{1:~b) a) $32$ \ b) $10\,000$ \ c) $7^3 = 343$ \ d) $3^5 = 243$}}
\end{pfaufg}

% L1-A: Potenz als Malkette
\begin{pfaufg}{}{2.}{}
Berechne beides. \ a)~$12^2$ und $12 \cdot 2$ \qquad b)~$2^6$ und $2 \cdot 6$ \qquad c)~$5^3$ und $5 \cdot 3$ \qquad d)~$15^2$ und $15 \cdot 2$\karo{3}
\fusshilfe{\mbox{2:~a) $144$}}
\end{pfaufg}

% L1-A: Potenz als Malkette
\begin{pfaufg}{}{3.}{}
Ein Bakterium teilt sich jede Stunde in zwei. \ a)~Wie viele Bakterien sind es nach 1, 2 und 3 Stunden? \ b)~Schreibe die Anzahl nach 10 Stunden als Potenz und rechne sie aus.\karo{3}
\fusshilfe{\mbox{3:~a) $2$; $4$; $8$ \ b) $2^{10} = 1024$}}
\end{pfaufg}

% L1-A: Potenz als Malkette
\begin{pfaufg}{}{4.}{}
Ein Blatt Papier ist $0{,}1$\,mm dick. Du faltest es in der Mitte, dann noch einmal, insgesamt 7-mal. Jedes Falten verdoppelt die Zahl der Lagen. Ist der Stapel dann dicker als 1\,cm?\karo{3}
\fusshilfe{\mbox{4:~ja: $2^7 = 128$ Lagen, $12{,}8$\,mm}}
\end{pfaufg}

% L1-B: Minus und Komma in der Potenz
\begin{pfaufg}{}{5.}{}
\grau{a)\ Berechne \ a)~$(-5)^2$ \qquad b)~$(-4)^3$ \qquad c)~$-5^2$ \qquad d)~$0{,}3^2$ \qquad e)~$\left(\frac{3}{4}\right)^2$\par\smallskip Klammer: Die ganze Zahl $-5$ wird malgenommen. Minus mal Minus gibt Plus: \quad $(-5)^2 = (-5) \cdot (-5) = 25$\par Drei Faktoren mit Minus: \quad $(-4) \cdot (-4) \cdot (-4) = 16 \cdot (-4) = -64$\par Keine Klammer: Der Exponent gehört nur zur $5$, das Minus bleibt davor: \quad $-5^2 = -(5 \cdot 5) = -25$\par Dezimalzahl: \quad $0{,}3 \cdot 0{,}3$; \ $3 \cdot 3 = 9$, Kommastellen $1 + 1 = 2$: \quad $0{,}3^2 = 0{,}09$\par Bruch: Zähler mal Zähler, Nenner mal Nenner: \quad $\left(\frac{3}{4}\right)^2 = \frac{3 \cdot 3}{4 \cdot 4} = \frac{9}{16}$}
\par\medskip
b)\ Berechne \ a)~$(-6)^2$ \qquad b)~$(-2)^5$ \qquad c)~$(-1)^8$ \qquad d)~$(-10)^3$\karo{3}
\fusshilfe{\mbox{5:~b) a) $36$ \ b) $-32$ \ c) $1$ \ d) $-1000$}}
\end{pfaufg}

% L1-B: Minus und Komma in der Potenz
\begin{pfaufg}{}{6.}{}
Berechne beides. \ a)~$(-7)^2$ und $-7^2$ \qquad b)~$(-1)^4$ und $-1^4$ \qquad c)~$(-3)^3$ und $-3^3$\karo{3}
\fusshilfe{\mbox{6:~a) $49$}}
\end{pfaufg}

% L1-B: Minus und Komma in der Potenz
\begin{pfaufg}{}{7.}{}
Berechne \ a)~$0{,}5^2$ \qquad b)~$1{,}2^2$ \qquad c)~$0{,}2^3$ \qquad d)~$\left(\frac{2}{5}\right)^2$ \qquad e)~$\left(\frac{1}{2}\right)^3$\karo{3}
\fusshilfe{\mbox{7:~a) $0{,}25$ \ b) $1{,}44$ \ c) $0{,}008$ \ d) $\frac{4}{25}$ \ e) $\frac{1}{8}$}}
\end{pfaufg}

% L1-B: Minus und Komma in der Potenz
\begin{pfaufg}{}{8.}{}
Hier stehen vier Zahlen. Unterstreiche die kleinste. \quad $-6^2$ \qquad $(-6)^2$ \qquad $(-3)^3$ \qquad $-0{,}5^2$\karo{3}
\fusshilfe{\mbox{8:~$-6^2$ (Wert $-36$)}}
\end{pfaufg}

% L1-C: Vergleichen und Exponent finden
\begin{pfaufg}{}{9.}{}
\grau{a)\ a)~Welche Zahl ist größer, $2^5$ oder $5^2$? \qquad b)~Für welches $x$ gilt $3^x = 243$?\par\smallskip Beide ausrechnen: \quad $2^5 = 2 \cdot 2 \cdot 2 \cdot 2 \cdot 2 = 32$, \quad $5^2 = 5 \cdot 5 = 25$. \ Also $2^5 > 5^2$.\par Exponent finden: Die $3$ immer wieder malnehmen, bis $243$ dasteht: \quad $3,\ 9,\ 27,\ 81,\ 243$\par Faktoren zählen: $243 = 3 \cdot 3 \cdot 3 \cdot 3 \cdot 3$, das sind fünf. \ Also $x = 5$.\par Probe: \quad $3^5 = 243$ \checkmark}
\par\medskip
b)\ Setze $<$, $=$ oder $>$ ein. \ a)~$3^2 \;\square\; 2^3$ \qquad b)~$10^2 \;\square\; 2^7$ \qquad c)~$4^2 \;\square\; 2^4$ \qquad d)~$1^{10} \;\square\; 10^1$\karo{3}
\fusshilfe{\mbox{9:~b) a) $>$ \ b) $<$ \ c) $=$ \ d) $<$}}
\end{pfaufg}

% L1-C: Vergleichen und Exponent finden
\begin{pfaufg}{}{10.}{}
Für welches $x$ gilt die Gleichung? \ a)~$2^x = 64$ \qquad b)~$5^x = 625$ \qquad c)~$10^x = 100\,000$ \qquad d)~$3^x = 27$\karo{3}
\fusshilfe{\mbox{10:~a) $x = 6$ \ b) $x = 4$ \ c) $x = 5$ \ d) $x = 3$}}
\end{pfaufg}

% L1-C: Vergleichen und Exponent finden
\begin{pfaufg}{}{11.}{}
Hier stehen drei Zahlen. Unterstreiche die größte. \quad $5^3$ \qquad $2^7$ \qquad $11^2$\karo{3}
\fusshilfe{\mbox{11:~$2^7$ (Wert $128$)}}
\end{pfaufg}

% L1-C: Vergleichen und Exponent finden
\begin{pfaufg}{}{12.}{}
Auf einem Teich wachsen Algen. Heute bedecken sie 1\,m$^2$, jeden Tag verdoppelt sich die Fläche. Der Teich ist 512\,m$^2$ groß. Ist er in einer Woche ganz zugewachsen?\karo{3}
\fusshilfe{\mbox{12:~nein: $2^9 = 512$, nach $9$ Tagen}}
\end{pfaufg}

% L1-D: Negative Hochzahl und hoch null
\begin{pfaufg}{}{13.}{}
\grau{a)\ Schreibe $5^{-2}$ als Bruch und als Dezimalzahl. Ist $5^{-2}$ größer oder kleiner als $0{,}05$? Berechne auch $9^0$.\par\smallskip Minus in der Hochzahl heißt „1 geteilt durch“: \quad $5^{-2} = \frac{1}{5^2}$\par Unten ausrechnen: $5^2 = 25$, \ also \ $5^{-2} = \frac{1}{25}$\par Als Dezimalzahl: erweitern auf Hundertstel: \quad $\frac{1}{25} = \frac{4}{100} = 0{,}04$ \quad (oder $1 : 25$)\par Vergleichen: \quad $0{,}04 < 0{,}05$, \ also \ $5^{-2} < 0{,}05$\par Hoch null: $9^2 = 81$, \ $9^1 = 9$, \ $9^0 = 9 : 9 = 1$ \ (jede Stufe tiefer: durch $9$ teilen).}
\par\medskip
b)\ Schreibe als Bruch. \ a)~$3^{-2}$ \qquad b)~$10^{-3}$ \qquad c)~$2^{-4}$ \qquad d)~$7^{-1}$ \qquad e)~$6^0$\karo{3}
\fusshilfe{\mbox{13:~b) a) $\frac{1}{9}$ \ b) $\frac{1}{1000}$ \ c) $\frac{1}{16}$ \ d) $\frac{1}{7}$ \ e) $1$}}
\end{pfaufg}

% L1-D: Negative Hochzahl und hoch null
\begin{pfaufg}{}{14.}{}
Schreibe als Bruch und als Dezimalzahl. \ a)~$2^{-1}$ \qquad b)~$5^{-1}$ \qquad c)~$10^{-2}$ \qquad d)~$4^{-1}$\karo{3}
\fusshilfe{\mbox{14:~a) $\frac{1}{2} = 0{,}5$ \ b) $\frac{1}{5} = 0{,}2$ \ c) $\frac{1}{100} = 0{,}01$ \ d) $\frac{1}{4} = 0{,}25$}}
\end{pfaufg}

% L1-D: Negative Hochzahl und hoch null
\begin{pfaufg}{}{15.}{}
Setze $<$, $=$ oder $>$ ein. \ a)~$8^{-1} \;\square\; 0{,}1$ \qquad b)~$10^{-3} \;\square\; 0{,}01$ \qquad c)~$2^{-2} \;\square\; 0{,}2$ \qquad d)~$3^0 \;\square\; 0{,}9$\karo{3}
\fusshilfe{\mbox{15:~a) $>$ \ b) $<$ \ c) $>$ \ d) $>$}}
\end{pfaufg}

% L1-D: Negative Hochzahl und hoch null
\begin{pfaufg}{}{16.}{}
Ordne die Zahlen der Größe nach. Beginne mit der kleinsten: \quad $5^{-1}$ \qquad $0{,}15$ \qquad $20^{-1}$ \qquad $0{,}3$ \qquad $4^{-1}$\karo{3}
\fusshilfe{\mbox{16:~$20^{-1} < 0{,}15 < 5^{-1} < 4^{-1} < 0{,}3$}}
\end{pfaufg}

% L1-T: Probetest
\begin{pfaufg}{}{17.}{}
Für welches $x$ gilt $8^x = 4096$?\karo{3}
\fusshilfe{\mbox{17:~$x = 4$}}
\end{pfaufg}

% L1-T: Probetest
\begin{pfaufg}{}{18.}{}
Berechne \ a)~$(-6)^3$ \qquad b)~$-6^2$ \qquad c)~$0{,}4^3$ \qquad d)~$(-5)^0$ \qquad e)~$10^{-1}$\karo{3}
\fusshilfe{\mbox{18:~a) $-216$ \ b) $-36$ \ c) $0{,}064$ \ d) $1$ \ e) $0{,}1$}}
\end{pfaufg}

% L1-T: Probetest
\begin{pfaufg}{}{19.}{}
Hier stehen drei Zahlen. Unterstreiche die größte. \quad $4^4$ \qquad $3^5$ \qquad $15^2$\karo{3}
\fusshilfe{\mbox{19:~$4^4$ (Wert $256$)}}
\end{pfaufg}

% L1-T: Probetest
\begin{pfaufg}{}{20.}{}
Setze $<$, $=$ oder $>$ ein. \qquad $2^{-4} \;\square\; 0{,}06$\karo{3}
\fusshilfe{\mbox{20:~$>$}}
\end{pfaufg}

% L1-T: Probetest
\begin{pfaufg}{}{21.}{}
Ein Video wird geteilt: Wer es bekommt, schickt es nach einer Stunde an 4 neue Personen. Am Anfang hat es eine Person. Lena sagt: „In der 5. Stunde bekommen es schon mehr als 1000 Personen neu.“ Stimmt das?\karo{3}
\fusshilfe{\mbox{21:~ja: $4^5 = 1024$}}
\end{pfaufg}

\par\vfill\noindent{\footnotesize Vorher: Rechnen mit negativen Zahlen\quad\textbullet\quad Weiter: Zehnerpotenzen\par}
\end{document}
